
\section{L'interfaccia unica del kernel locale}
\label{sec:interfaccia}

Ogni backend implementa tre operazioni fondamentali per la gestione e l'esecuzione del kernel locale, oltre ad alcune funzioni ausiliarie utilizzate per la raccolta delle metriche sperimentali.

\begin{lstlisting}[language=C,caption={Interfaccia di \file{src/kernel/kernel.h}.}]
/* esecuzione */
local_gemm_t *local_gemm_create(int m, int n, int k,
                                const scalar_t *A_loc,
                                int lda, int ldx, int ldy);
void local_gemm(local_gemm_t *ctx,
                const scalar_t *RESTRICT X, int ldx,
                scalar_t *RESTRICT Y, int ldy);
void local_gemm_destroy(local_gemm_t *ctx);

\end{lstlisting}
\color{red}
L'interfaccia segue un ciclo di vita composto da tre fasi rappresentate da altrettante funzioni che definiscono il ciclo di vita del kernel locale:
\code{local\_gemm\_create} ne inizializza il contesto,
\code{local\_gemm} esegue il calcolo e
\code{local\_gemm\_destroy} rilascia le risorse associate.
\color{black}


Inoltre l'interfaccia espone alcune funzioni di introspezione utilizzate dal driver di
benchmark per raccogliere separatamente tempi di calcolo, trasferimenti, costi di
inizializzazione e parametri di esecuzione del backend. Tali metriche vengono descritte nella sezione dedicata alla metodologia sperimentale.

\subsection{Ciclo di vita del kernel}
\label{sec:tre-fasi}

La separazione tra \code{local\_gemm\_create}, \code{local\_gemm} e
\code{local\_gemm\_destroy} deriva dal diverso ciclo di vita dei dati
e delle risorse coinvolte nel calcolo.

La matrice $A$ rimane invariata durante tutte le ripetizioni, mentre
$X$ e $Y$ sono coinvolte in ogni esecuzione del prodotto locale.
Questa differenza permette di separare le operazioni eseguite una sola
volta da quelle necessarie a ogni invocazione.

\begin{enumerate}

    \item \textbf{Inizializzazione (\code{local\_gemm\_create}).}
    Viene creato il contesto del backend e vengono predisposte le
    risorse riutilizzabili durante le successive esecuzioni.
    
    Nei backend CPU vengono memorizzate le informazioni necessarie
    all'esecuzione, tra cui le dimensioni del problema e il riferimento
    al blocco locale di $A$.
    
    Nei backend CUDA vengono inoltre allocati i buffer sul device e
    $A$ viene trasferita una sola volta nella memoria della GPU.

    \item \textbf{Esecuzione (\code{local\_gemm}).}
    Questa fase viene ripetuta per ogni invocazione del prodotto
    locale. La matrice $X$, precedentemente distribuita tra i processi tramite
    \code{MPI\_Bcast}, viene utilizzata insieme al blocco locale di $A$
    per produrre il contributo locale a $Y$.
    
    Nei backend CUDA vengono riutilizzati i buffer allocati durante
    l'inizializzazione: la matrice $X$ viene trasferita sul device, viene eseguito
    il kernel e il risultato $Y$ viene riportato in memoria host.
    
    I contributi prodotti dai processi vengono successivamente
    combinati tramite \code{MPI\_Reduce}.

    \item \textbf{Finalizzazione (\code{local\_gemm\_destroy}).}
    Al termine delle ripetizioni vengono rilasciate le risorse
    associate al contesto, comprese le eventuali allocazioni sulla GPU.

\end{enumerate}

Questa organizzazione evita di ripetere a ogni invocazione le
operazioni associate ad $A$ e separa il costo di inizializzazione
dal costo delle successive esecuzioni del prodotto locale.

\subsection{Contesto opaco del backend}

Lo stato specifico di ciascun backend è mantenuto all'interno del 
tipo opaco \code{local\_gemm\_t}. La struttura concreta è definita
separatamente da ogni implementazione e può quindi contenere
informazioni differenti senza modificare l'interfaccia comune.

Ad esempio, un backend CPU può memorizzare le dimensioni del problema
e il riferimento ad $A$, mentre un backend CUDA mantiene anche i
puntatori alla memoria del device e le risorse necessarie
all'esecuzione. I dettagli interni del backend rimangono così
nascosti al resto dell'applicazione.

\subsection{Configurazione delle build}

Le principali scelte sperimentali, tra cui backend di calcolo,
precisione numerica e alcuni parametri dei kernel CUDA, vengono
selezionate a tempo di compilazione tramite il \code{Makefile}.

Le principali varianti di configurazione utilizzate nelle campagne
sperimentali vengono associate a binari distinti, evitando che build
differenti si sovrascrivano e rendendo immediatamente identificabile
la configurazione utilizzata.